%%% =======================================================================
%%% ISLS Extension Phase 10 v1.0.0 (Revised)
%%% C28: Babylon Bridge — Compiler-Backed Multi-Language Forge
%%% Author: Sebastian Klemm
%%% Date: 18 March 2026
%%% Status: Implementation Specification for Claude Code
%%% =======================================================================
\documentclass[11pt,a4paper,twoside]{article}

\usepackage[utf8]{inputenc}
\usepackage{amsmath,amssymb,amsthm}
\usepackage{algorithm,algorithmic}
\usepackage{booktabs,longtable,tabularx}
\usepackage{enumitem}
\usepackage{geometry}
\usepackage{hyperref}
\usepackage{cleveref}
\usepackage{xcolor}
\usepackage{fancyhdr}
\usepackage{listings}
\usepackage{titlesec}

\geometry{margin=2.5cm,top=3cm,bottom=3cm,headheight=14pt}

\theoremstyle{definition}
\newtheorem{definition}{Definition}[section]
\newtheorem{axiom}{Axiom}[section]
\newtheorem{requirement}{Requirement}[section]
\newtheorem{invariant}{Invariant}[section]

\theoremstyle{plain}
\newtheorem{theorem}{Theorem}[section]

\theoremstyle{remark}
\newtheorem{remark}{Remark}[section]

\newcommand{\ISLS}{\textsf{ISLS}}

\definecolor{sec-color}{HTML}{1e3a5f}

\titleformat{\section}
  {\Large\bfseries\color{sec-color}}
  {\thesection}{1em}{}

\lstset{
  basicstyle=\ttfamily\small,
  breaklines=true,
  frame=single,
  backgroundcolor=\color{gray!5},
  rulecolor=\color{gray!30},
}

\pagestyle{fancy}
\fancyhf{}
\fancyhead[LE,RO]{\ISLS{} Phase 10 v1.0.0 (Revised)}
\fancyhead[RE,LO]{Sebastian Klemm}
\fancyfoot[C]{\thepage}

\hypersetup{
  colorlinks=true,
  linkcolor=blue!60!black,
  pdftitle={ISLS Phase 10 — Babylon Bridge},
  pdfauthor={Sebastian Klemm},
}

\begin{document}

\begin{titlepage}
\centering
\vspace*{2.5cm}
{\Huge\bfseries \ISLS{} Extension Phase~10\\[0.3cm]
v1.0.0 (Revised)}\\[1.5cm]
{\Large C28: The Babylon Bridge}\\[1cm]
{\large Compiler-Backed Multi-Language Forge\\[0.3cm]
Not ``more prompt templates.''\\
A real IR with real parsers, real embedding,\\
and real cross-language transformation.}\\[2cm]
{\large Sebastian Klemm}\\[0.5cm]
{\large 18 March 2026}\\[1.5cm]

\begin{abstract}
\noindent
This document supersedes the earlier Phase~10 specification. Instead
of adding six LLM-prompt matrices for different languages (hoping the
Oracle outputs valid TypeScript, Python, SQL), this revision
integrates the Babylon Compiler as a structural backbone.

Babylon brings: three working parsers (Sanskroot, HanLan, Cuneiform)
that lower to a shared \texttt{glyph-ir::IrDocument}; a 5D H5
embedding engine; Q16 fixed-point determinism; a crystallization and
gating pipeline; and a verified bundle emitter. Merkaba Factory brings:
codegen backends for Python, Rust, JavaScript, TypeScript, and Go.

The bridge: the \ISLS{} Forge generates an \texttt{ArtifactIR}
(C22). The Babylon Bridge translates ArtifactIR components into
\texttt{glyph-ir::IrDocument} fragments. The Oracle (C25) fills
semantic gaps (function bodies, logic). The filled IR is validated by
Babylon's pipeline (embedding, gating, evidence chain). Then codegen
backends emit to the target language.

The result: generated code is not a string the LLM produced. It is
an IR that was parsed, embedded, gated, and emitted through a real
compiler pipeline. If the code has structural errors, the IR
validation catches them \emph{before} the toolchain runs. If the code
needs to be in a different language, re-emit from the same IR.

\medskip\noindent\textbf{Status:} Implementation Specification.\\
\textbf{Absorbs:} Babylon Compiler (\texttt{glyph-ir},
  \texttt{glyph-embed}, \texttt{glyph-gate}, \texttt{glyph-mef},
  \texttt{glyph-frontends}); Merkaba \texttt{domain-translator}
  codegen concept.\\
\textbf{Depends on:} C22 (ArtifactIR), C23 (Forge), C25 (Oracle),
  C27 (Foundry).
\end{abstract}
\end{titlepage}

\tableofcontents
\newpage


%%% =====================================================================
\part{Why This Is Different}\label{part:why}

\section{The Previous Approach (Discarded)}\label{sec:old}

The original Phase~10 added six ``Matrix'' implementations: each was
a prompt template that told the Oracle ``generate TypeScript'' or
``generate Python.'' The output was a string. Validation was: run
\texttt{tsc} or \texttt{python -m py\_compile} and hope.

Problems:
\begin{enumerate}[nosep]
  \item \textbf{No structural validation}: a syntactically valid
        TypeScript file can be architecturally wrong (wrong imports,
        wrong types, missing interfaces). The toolchain check catches
        syntax, not structure.
  \item \textbf{No cross-language consistency}: if the Rust backend
        defines a \texttt{User} struct, the TypeScript frontend must
        have an identical \texttt{User} interface. LLM-only generation
        has no mechanism to enforce this.
  \item \textbf{No pivot representation}: to support $n$ languages, you
        need $n$ prompt templates. To support $n \times m$
        cross-language transforms, you need $n \times m$ templates.
        With an IR, you need $n$ frontends + $m$ backends.
\end{enumerate}

\section{The Babylon Approach}\label{sec:new}

\begin{lstlisting}[title={Previous: String-based}]
Oracle("generate Rust code") -> String -> cargo check -> hope
Oracle("generate TypeScript") -> String -> tsc --noEmit -> hope
\end{lstlisting}

\begin{lstlisting}[title={Revised: IR-based}]
ArtifactIR
  -> Babylon Bridge: convert to glyph-ir IrDocument
  -> Oracle fills function bodies as IR node content
  -> glyph-embed: compute H5 embedding (structural validation)
  -> glyph-gate: PoR check (structural consistency)
  -> Codegen Backend: emit to target language from validated IR
  -> Toolchain: cargo check / tsc / python (final confirmation)
\end{lstlisting}

The IR is the single source of truth. Languages are views of the IR,
not independent LLM outputs.


%%% =====================================================================
\part{Architecture}\label{part:arch}

\section{The Bridge Pipeline}\label{sec:pipeline}

\begin{lstlisting}[title={Complete Pipeline}]
DecisionSpec
  |
  v
[PMHD Drill] -> PmhdMonolith
  |
  v
[ArtifactIR] -> Components, Interfaces, Constraints
  |
  v
[Babylon Bridge: ArtifactIR -> glyph-ir]
  |                                       (1) Structure from ArtifactIR
  |                                       (2) Bodies from Oracle
  v
[glyph-ir::IrDocument] -- canonical IR with nodes, edges, embedding
  |
  v
[glyph-embed] -- H5 embedding: structural_coupling, functional_density,
  |               topological_complexity, symmetry, entropy
  v
[glyph-gate] -- PoR check: does the IR represent a coherent program?
  |
  v
[glyph-mef] -- Evidence chain for the compilation
  |
  v
[Codegen Backend] -- Emit from IR to target language
  |    |    |    |    |
  v    v    v    v    v
 Rust  TS  Python SQL  YAML
  |
  v
[Foundry] -- cargo check / tsc / pytest (final toolchain validation)
  |
  v
[Crystal] -- proven to: (a) have valid IR structure,
             (b) pass H5 embedding checks,
             (c) pass PoR gate,
             (d) compile in target language,
             (e) pass tests
\end{lstlisting}


\section{What Babylon Provides}\label{sec:babylonprovides}

\begin{tabularx}{\textwidth}{l l X}
\toprule
\textbf{Crate} & \textbf{Used For} & \textbf{Integration Point} \\
\midrule
\texttt{glyph-ir} & Universal IR & ArtifactIR $\to$ IrDocument
  conversion. Nodes: Module, Function, Block, If, Return, Call,
  Assignment, Literal, BinaryOp, Identifier. Edges: Contains, Next,
  Condition, ThenBranch, ElseBranch, Argument, Target, Value, Left,
  Right, CalleeRef. \\
\texttt{glyph-embed} & 5D structural validation & After IR
  construction, compute H5 embedding. Five axes: structural coupling,
  functional density, topological complexity, symmetry, entropy.
  Embedding scores below threshold flag structural problems. \\
\texttt{glyph-canon} & Canonical serialization & JCS (RFC 8785) +
  SHA-256 digest. Same canonical form as \ISLS{}. \\
\texttt{glyph-q16} & Fixed-point arithmetic & Q16.16 for embedding
  axes. Deterministic across platforms. \\
\texttt{glyph-gate} & Structural gating & PoR FSM + obligation
  evaluation on the IR. Rejects structurally inconsistent programs
  before codegen. \\
\texttt{glyph-mef} & Evidence chain & Append-only hash chain for
  compilation steps. Fuses with \ISLS{} evidence chain. \\
\texttt{glyph-frontends} & (Future) Parse LLM output & When the
  Oracle produces code in a language that Babylon can parse, the
  output is parsed back to IR for round-trip validation. \\
\bottomrule
\end{tabularx}


%%% =====================================================================
\part{The Bridge: ArtifactIR $\to$ glyph-ir}\label{part:bridge}

\section{Conversion}\label{sec:convert}

\begin{definition}[ArtifactIR to IrDocument]\label{def:convert}
Each ArtifactIR Component maps to glyph-ir nodes:
\begin{lstlisting}
pub fn artifact_to_ir(ir: &ArtifactIR) -> IrDocument {
    let mut doc = IrDocument::new(
        &ir.header.domain,
        &hex_encode(&ir.header.artifact_id),
    );

    // Root module node
    let root_id = format!("n_root");
    doc.nodes.push(IrNode::new(&root_id, NodeKind::Module, "root"));

    // Each ArtifactIR component -> Function node
    for (i, comp) in ir.components.iter().enumerate() {
        let fn_id = format!("n_fn_{}", i);
        let mut node = IrNode::new(&fn_id, NodeKind::Function, &comp.name);
        node.properties = Some({
            let mut props = BTreeMap::new();
            props.insert("kind".into(), json!(comp.kind));
            props.insert("content".into(), json!(comp.content));
            if !comp.dependencies.is_empty() {
                props.insert("deps".into(), json!(comp.dependencies));
            }
            props
        });
        doc.nodes.push(node);
        doc.edges.push(IrEdge::new(&root_id, &fn_id, EdgeKind::Contains));
    }

    // Each ArtifactIR interface -> edges between function nodes
    for iface in &ir.interfaces {
        if let (Some(p), Some(c)) = (
            find_node_idx(&doc, &iface.provider),
            find_node_idx(&doc, &iface.consumer),
        ) {
            doc.edges.push(IrEdge::new(
                &doc.nodes[p].id,
                &doc.nodes[c].id,
                EdgeKind::CalleeRef,
            ));
        }
    }

    doc.canonicalize();
    doc
}
\end{lstlisting}
\end{definition}

\begin{definition}[Oracle-Filled IR]\label{def:oraclefill}
After conversion, each Function node has a \texttt{content} property
containing the ArtifactIR component description. The Oracle (C25)
is called to fill each node with implementation code:
\begin{lstlisting}
pub fn fill_ir_with_oracle(
    doc: &mut IrDocument,
    oracle: &mut OracleEngine,
    domain: &str,
) -> Result<()> {
    for node in &mut doc.nodes {
        if node.kind == NodeKind::Function {
            let content = node.properties
                .as_ref()
                .and_then(|p| p.get("content"))
                .and_then(|v| v.as_str())
                .unwrap_or("");

            let prompt = format!(
                "Implement this {} function:\n{}\n\
                 Output ONLY the function body code. No explanation.",
                domain, content
            );

            let response = oracle.synthesize_raw(&prompt)?;

            // Store the implementation in the node
            if let Some(props) = &mut node.properties {
                props.insert("implementation".into(),
                    json!(response.content));
            }
        }
    }
    Ok(())
}
\end{lstlisting}
The key difference: the Oracle fills \emph{nodes in an IR}, not
\emph{entire files}. The structure (which functions exist, how they
connect, what their signatures are) comes from the ArtifactIR. Only
the bodies come from the Oracle.
\end{definition}


\section{H5 Embedding as Structural Gate}\label{sec:h5}

\begin{definition}[Structural Validation via Embedding]\label{def:h5val}
After the Oracle fills the IR, the H5 embedding is computed:
\begin{lstlisting}
let embedding = glyph_embed::compute_embedding(&doc);
\end{lstlisting}

The five axes catch structural problems:
\begin{itemize}[nosep]
  \item $a_1$ (structural coupling): if too high, components are
        over-coupled. If zero, components are disconnected.
  \item $a_2$ (functional density): if zero, no functions exist. If
        too high, too many tiny functions (fragmentation).
  \item $a_3$ (topological complexity): if too high, too many
        branches/conditions (oracle over-complicated it).
  \item $a_4$ (symmetry): if too low, asymmetric structure (some
        modules much larger than others).
  \item $a_5$ (entropy): if too low, repetitive code. If too high,
        random/incoherent.
\end{itemize}

Each axis has a range check. If any axis is outside $[\theta_{\min},
\theta_{\max}]$, the IR is flagged for re-generation.
\end{definition}


%%% =====================================================================
\part{Codegen Backends}\label{part:codegen}

\section{Backend Trait}\label{sec:backendtrait}

\begin{definition}[Codegen Backend]\label{def:backend}
\begin{lstlisting}
pub trait CodegenBackend: Send + Sync {
    /// Target language name
    fn language(&self) -> &str;

    /// File extension
    fn extension(&self) -> &str;

    /// Emit a complete source file from an IrDocument
    fn emit(&self, doc: &IrDocument, config: &CodegenConfig) -> Result<Vec<EmittedFile>>;

    /// Emit a single function node to a code string
    fn emit_function(&self, node: &IrNode, doc: &IrDocument) -> Result<String>;
}

pub struct EmittedFile {
    pub path: String,      // relative path (e.g., "src/main.rs")
    pub content: String,
    pub language: String,
}

pub struct CodegenConfig {
    pub indent: usize,
    pub line_width: usize,
    pub include_tests: bool,
    pub include_docs: bool,
}
\end{lstlisting}
\end{definition}


\section{Built-In Backends}\label{sec:backends}

\begin{tabularx}{\textwidth}{l l l X}
\toprule
\textbf{Backend} & \textbf{Language} & \textbf{Ext} & \textbf{What It Emits} \\
\midrule
\texttt{RustBackend} & Rust & .rs & \texttt{mod}, \texttt{pub fn},
  \texttt{struct}, \texttt{impl}, \texttt{\#[test]}. Cargo.toml
  generation. \\
\texttt{TypeScriptBackend} & TypeScript & .ts/.tsx & \texttt{export
  function}, \texttt{interface}, \texttt{type}, React components
  (if .tsx). \\
\texttt{PythonBackend} & Python & .py & \texttt{def}, \texttt{class},
  \texttt{@dataclass}, type hints, \texttt{pytest} tests. \\
\texttt{SqlBackend} & SQL & .sql & \texttt{CREATE TABLE},
  \texttt{ALTER TABLE}, indexes, views, migration files. \\
\texttt{YamlBackend} & YAML & .yaml/.yml & Docker Compose,
  GitHub Actions, Kubernetes manifests. \\
\texttt{MarkdownBackend} & Markdown & .md & README, API docs,
  architecture docs, generated from IR structure. \\
\bottomrule
\end{tabularx}

\begin{definition}[Multi-Target Emission]\label{def:multitarget}
A single IR can be emitted to multiple targets simultaneously:
\begin{lstlisting}
pub fn emit_multi(
    doc: &IrDocument,
    targets: &[(&dyn CodegenBackend, CodegenConfig)],
) -> Result<Vec<EmittedFile>> {
    let mut all_files = Vec::new();
    for (backend, config) in targets {
        all_files.extend(backend.emit(doc, config)?);
    }
    Ok(all_files)
}
\end{lstlisting}
This is how a full-stack project is produced: one IR, multiple
backends, one consistent set of types.
\end{definition}


\section{Cross-Language Type Consistency}\label{sec:crosslang}

\begin{definition}[Type Bridge]\label{def:typebridge}
Types defined in one language view of the IR are automatically
available in all other language views:
\begin{lstlisting}
// In the IR: a Module node "User" with properties {fields: [...]}
// Rust backend emits:
pub struct User {
    pub id: i64,
    pub name: String,
    pub email: String,
}

// TypeScript backend emits (from the same IR node):
export interface User {
    id: number;
    name: string;
    email: string;
}

// Python backend emits:
@dataclass
class User:
    id: int
    name: str
    email: str

// SQL backend emits:
CREATE TABLE users (
    id   BIGINT PRIMARY KEY,
    name TEXT NOT NULL,
    email TEXT NOT NULL
);
\end{lstlisting}

The types are consistent because they come from the same IR node. No
``hoping the LLM generates matching types.''
\end{definition}


%%% =====================================================================
\part{Round-Trip Validation}\label{part:roundtrip}

\section{Parse-Back Validation}\label{sec:parseback}

\begin{definition}[Round-Trip Check]\label{def:roundtrip}
For languages where Babylon has a frontend parser (currently:
Sanskroot, HanLan, Cuneiform; extendable to Rust, Python, TS), the
emitted code is parsed back to IR and compared:
\begin{lstlisting}
pub fn round_trip_validate(
    original_ir: &IrDocument,
    emitted_code: &str,
    frontend: &str,
) -> Result<RoundTripResult> {
    let reparsed_ir = parse_to_ir(emitted_code, frontend)?;
    let original_embedding = compute_embedding(original_ir);
    let reparsed_embedding = compute_embedding(&reparsed_ir);
    let distance = h5_distance(&original_embedding, &reparsed_embedding);

    Ok(RoundTripResult {
        parseable: true,
        embedding_distance: distance,
        structurally_equivalent: distance < ROUND_TRIP_THRESHOLD,
    })
}
\end{lstlisting}

If the round-trip distance exceeds the threshold, the codegen backend
has a bug or the Oracle's body code introduced structural changes that
broke the IR's topology. This is caught at the IR level, not at the
toolchain level.
\end{definition}


%%% =====================================================================
\part{Full-Stack Templates (IR-Based)}\label{part:templates}

\section{IR-Based Templates}\label{sec:irtemplates}

The six full-stack templates from the original Phase~10 (T11--T16) are
reimplemented as IR trees rather than prompt templates. Each template
is a pre-built \texttt{IrDocument} with structural nodes and edges,
ready for Oracle body-filling and multi-target emission.

\begin{tabularx}{\textwidth}{l l r X}
\toprule
\textbf{ID} & \textbf{Name} & \textbf{IR Nodes} &
  \textbf{Emit Targets} \\
\midrule
T11 & SaaS Starter & 45 & Rust + TS + SQL + Docker + MD \\
T12 & Dashboard & 32 & Rust + TS + SQL + Docker \\
T13 & API + Docs & 20 & Rust + MD + Docker \\
T14 & ML Pipeline & 18 & Python + SQL + Docker \\
T15 & Static Site & 14 & HTML/CSS + MD \\
T16 & Monorepo & 38 & Rust + TS + SQL + Docker + MD \\
\bottomrule
\end{tabularx}

Each template's IR is pre-embedded (H5 axes computed), pre-gated
(PoR validated), and stored as a crystal in the template catalog.
The Oracle fills only the function bodies. Structure is guaranteed
by the IR.


%%% =====================================================================
\part{Integration with Existing Crates}\label{part:integration}

\section{Dependency Strategy}\label{sec:deps}

The Babylon crates (\texttt{glyph-ir}, \texttt{glyph-embed},
\texttt{glyph-canon}, \texttt{glyph-q16}, \texttt{glyph-gate},
\texttt{glyph-mef}) are added as workspace dependencies. They are
\textbf{not} copied into \ISLS{} --- they are used as external crates
from the babylon-compiler repository (local path or git dependency).

\begin{lstlisting}[title={Workspace Cargo.toml additions}]
[workspace.dependencies]
glyph-ir     = { path = "../../babylon-compiler/crates/glyph-ir" }
glyph-embed  = { path = "../../babylon-compiler/crates/glyph-embed" }
glyph-canon  = { path = "../../babylon-compiler/crates/glyph-canon" }
glyph-q16    = { path = "../../babylon-compiler/crates/glyph-q16" }
glyph-gate   = { path = "../../babylon-compiler/crates/glyph-gate" }
glyph-mef    = { path = "../../babylon-compiler/crates/glyph-mef" }
\end{lstlisting}


\section{Modified Forge Pipeline}\label{sec:modforge}

\begin{lstlisting}[title={ForgeEngine with Babylon Bridge}]
impl ForgeEngine {
    pub fn forge(&mut self, spec: DecisionSpec) -> Result<ForgeResult> {
        // 1. PMHD Drill -> Monolith (unchanged)
        let monolith = self.drill(&spec)?;

        // 2. ArtifactIR (unchanged)
        let artifact_ir = ArtifactIR::build_from_monolith(&monolith, &spec, 0)?;

        // 3. NEW: Convert to glyph-ir
        let mut ir_doc = babylon_bridge::artifact_to_ir(&artifact_ir);

        // 4. NEW: Oracle fills function bodies in the IR
        babylon_bridge::fill_ir_with_oracle(
            &mut ir_doc, &mut self.oracle, &spec.domain
        )?;

        // 5. NEW: H5 embedding (structural validation)
        let embedding = glyph_embed::compute_embedding(&ir_doc);
        if !babylon_bridge::embedding_in_range(&embedding) {
            // Structural problem detected -> re-generate
            return Err(ForgeError::StructuralValidation(embedding));
        }

        // 6. NEW: Codegen to target language(s)
        let backends = self.get_backends(&spec.domain)?;
        let files = babylon_bridge::emit_multi(&ir_doc, &backends)?;

        // 7. Write files + Foundry loop (unchanged)
        let result = self.foundry.fabricate(files, &spec)?;

        // 8. Crystallize (now includes IR digest + H5 embedding)
        self.crystallize(result, &ir_doc, &embedding)
    }
}
\end{lstlisting}


\section{CLI Extension}\label{sec:cli}

\begin{lstlisting}
# Forge with target language (uses Babylon backend)
isls forge --spec intent.json --domain rust
isls forge --spec intent.json --domain typescript
isls forge --spec intent.json --domain python

# Forge full-stack (multiple backends from one IR)
isls forge --spec intent.json --template saas-starter

# Inspect the glyph-ir for a forge result
isls forge --spec intent.json --domain rust --emit-ir ir.json

# Compute H5 embedding of existing source code
isls babylon embed --source src/main.rs --frontend rust
\end{lstlisting}


%%% =====================================================================
\part{Acceptance Tests}\label{part:tests}

\begin{longtable}{l l p{5.5cm}}
\toprule
\textbf{ID} & \textbf{Name} & \textbf{Procedure} \\
\midrule
\endfirsthead \midrule \endhead \bottomrule \endfoot
AT-BB1 & IR conversion & Convert ArtifactIR to IrDocument; verify
  node and edge counts match component and interface counts. \\
AT-BB2 & IR determinism & Convert same ArtifactIR twice; verify
  identical IrDocument digest. \\
AT-BB3 & H5 embedding & Compute embedding of IR with 5 functions;
  verify all 5 axes in valid Q16 range. \\
AT-BB4 & Rust codegen & Emit IrDocument via RustBackend; verify
  .rs files produced with correct function signatures. \\
AT-BB5 & TypeScript codegen & Emit via TypeScriptBackend; verify
  .ts files with correct interface definitions. \\
AT-BB6 & Python codegen & Emit via PythonBackend; verify .py files
  with correct class definitions. \\
AT-BB7 & SQL codegen & Emit via SqlBackend; verify .sql files with
  correct CREATE TABLE statements. \\
AT-BB8 & Cross-language types & Define a type in IR; emit to Rust
  and TypeScript; verify field names and types match. \\
AT-BB9 & Multi-target emission & Emit single IR to 3 backends;
  verify files for all 3 languages produced. \\
AT-BB10 & Structural rejection & Create IR with zero functions
  (a2=0); verify H5 gate rejects. \\
AT-BB11 & Full-stack forge & Forge SaaS Starter (T11); verify
  backend/, frontend/, database/ dirs with correct language files. \\
AT-BB12 & Oracle fills bodies & Create IR with 3 function nodes;
  fill with mock Oracle; verify all 3 have implementation property. \\
\end{longtable}


%%% =====================================================================
\part{Handoff}\label{part:handoff}

\section{Test Count}

\begin{tabular}{l r r}
\toprule
\textbf{Phase} & \textbf{New Tests} & \textbf{Cumulative} \\
\midrule
Phase 1--9 & 326 & 326 \\
Phase 10 revised (C28) & 12 & $\ge$ 338 \\
\bottomrule
\end{tabular}

\section{Completion Criteria}

\begin{enumerate}[nosep]
  \item ArtifactIR converts to glyph-ir deterministically.
  \item Oracle fills function bodies in IR nodes (not whole files).
  \item H5 embedding validates structural coherence before codegen.
  \item Six codegen backends produce valid target-language files.
  \item Cross-language type consistency from single IR source.
  \item Full-stack templates emit multi-language projects.
  \item Babylon crates integrated as workspace dependencies.
  \item Foundry validates emitted code with language-appropriate
        toolchains.
  \item Total tests $\ge 338$, zero warnings.
  \item \textbf{The IR is the truth. Languages are views.}
\end{enumerate}

\vfill
\begin{center}
\rule{0.5\textwidth}{0.4pt}\\[0.5cm]
{\small Generated for \ISLS{} Extension Phase 10 (Revised) v1.0.0.\\
Not ``more prompts.'' A real compiler.\\
Structure from the IR. Bodies from the Oracle.\\
Types consistent across languages by construction.\\
Validated before emission. Verified after.}
\end{center}

\end{document}
